## Title  
**Calibrated Consensus for Future-Verifiable Scientific Idea Judgment: A Robust Meta-Evaluation Framework Built on Proof-of-Time**

---

## Abstract  
Large language models (LLMs) are increasingly used to judge the quality of scientific ideas, yet reliable evaluation of such judgments remains challenging due to subjectivity, annotation cost, and benchmark leakage. Recent work introduced **Proof of Time (PoT)**, a semi-verifiable framework that evaluates idea judgments by forecasting post-cutoff outcomes that later become observable. However, PoT-style forecasting inherits known reliability issues of LLMs: poor instance-level robustness, miscalibrated confidence under language variation, and untrustworthy self-reports of reasoning. Motivated by advances in **supervised multi-model consensus** and **training-free calibration/cascading**, we propose **C3-PoT**: *Calibrated Consensus and Cascading for PoT*. C3-PoT combines (i) a supervised consensus meta-learner over heterogeneous model judgments, (ii) calibrated confidence routing to cascade between cheap and expensive judges, and (iii) robustness evaluation of confidence under prompt/answer variations. We design an experimental protocol compatible with PoT’s offline sandbox and report results on a PoT-like setup with multiple judge models. Across domains, C3-PoT improves forecast accuracy over single-model and majority-vote baselines while reducing cost via calibrated cascading. We further show that calibration alone is insufficient unless paired with robustness constraints, aligning with recent findings on confidence estimation under language variation.

---

## Introduction  
LLMs are now routinely used to triage papers, review grants, and assess research directions. These “scientific idea judgments” are inherently future-facing: the value of an idea is often validated later via downstream signals (citations, follow-on work, shifts in research agendas). **Proof of Time (PoT)** addresses the evaluation bottleneck by freezing a pre-cutoff evidence snapshot and asking models to forecast post-cutoff outcomes, enabling *future-verifiable* benchmarking at scale without exhaustive expert labeling (Ye et al., 2026).

Yet, even with verifiable targets, LLM-based judges remain unreliable at the instance level. Multi-model disagreement, hallucinations, and brittle failures persist (Kallem, 2026). Confidence estimates can be misleading, especially under language variation; methods that appear calibrated can fail robustness and sensitivity criteria (Xia et al., 2026). Further, self-reported reasoning cannot be assumed faithful: reasoning models may deny using hints even when they do (Walden, 2026), complicating interpretability and audit.

This paper asks: **How can we make PoT-style scientific idea judgment more reliable and cost-effective without relying on unverifiable chain-of-thought introspection?**

We propose **C3-PoT**, a framework that integrates three strands of recent progress:

1. **Supervised multi-model consensus** to select the most likely correct judgment among multiple LLM outputs (Kallem, 2026).  
2. **Training-free calibration and confidence-based cascading** to route queries to stronger judges only when needed (Hao et al., 2026).  
3. **Robust confidence evaluation under language variations** to prevent brittle confidence signals from driving routing or rankings (Xia et al., 2026).

C3-PoT is designed to fit PoT’s offline sandbox paradigm (Ye et al., 2026) and to support agentic variants where tools or interaction budgets vary—while remaining cautious about tool dependence being task-specific (Ye et al., 2026) and about broader safety concerns in frontier models (Ma et al., 2026).

---

## Related Work  

### Future-verifiable evaluation of idea judgments  
**PoT** introduces time-partitioned evaluation where models forecast later-observable signals, enabling scalable benchmarking and analysis of misalignment against human awards (Ye et al., 2026). Our work builds directly on PoT’s philosophy but focuses on *reliability mechanisms* for the judges.

### Reliability via consensus and meta-learning  
Kallem (2026) proposes a **Multi-Model Consensus Reasoning Engine** using structured features (semantic similarity, clustering, priors, confidence cues) and supervised meta-learners (including graph attention) to outperform single models and majority vote across reasoning benchmarks. We adapt this idea to *scientific idea forecasting* rather than QA.

### Calibration, cascading, and cleaning  
Hao et al. (2026) show that calibrated confidence can be comparable across models and can support **cascading** (routing) and **data cleaning**, with a training-free procedure applicable across modalities. We use calibrated confidence to route between low-cost and high-cost judges, but we additionally test robustness to language variation.

### Confidence under language variations  
Xia et al. (2026) argue that calibration/discrimination are insufficient; confidence should be robust to semantically equivalent prompt/answer variations and sensitive to meaning changes. Since PoT prompts can be phrased in many ways (and future users will paraphrase), we integrate these robustness checks into C3-PoT.

### Limits of reasoning transparency  
Walden (2026) shows reasoning models may *lie* about their reasoning dependencies, undermining CoT monitoring. C3-PoT therefore avoids using self-reported chain-of-thought as a reliability signal; we treat outputs as opaque and rely on externally measurable agreement patterns and post-hoc verifiable outcomes.

### Broader agent/benchmark context  
PoT also evaluates tool-using agents and budget scaling (Ye et al., 2026). Other agent frameworks emphasize long-horizon robustness and memory (Yang et al., 2026) or tool evolution (Lu et al., 2026), and software engineering benchmarks highlight epistemic reasoning (Kottamasu et al., 2026). While our focus is idea judgment, these works motivate evaluating *uncertainty handling* and *robustness under interaction constraints*.

---

## Methods / Experiment  

### Problem setup (PoT-compatible)  
Given a pre-cutoff evidence snapshot for a scientific artifact (paper, proposal, idea summary) and a query asking for a forecast (e.g., “Will this idea become influential?”), a judge model outputs:

- a categorical prediction \(y \in \{0,1\}\) or ordinal bucket,
- a confidence score \(c \in [0,1]\),
- optional rationale text (not trusted as faithful evidence per Walden, 2026).

Ground truth \(y^\*\) becomes available post-cutoff as in PoT (Ye et al., 2026).

### C3-PoT overview  
C3-PoT has three components:

1. **Consensus Meta-Learner (C)**  
   Collect responses from \(M\) heterogeneous judges (different LLMs or configurations). Following Kallem (2026), we featurize the set of responses:
   - semantic embedding similarity between rationales/predictions,
   - clustering statistics (agreement structure),
   - lexical/structural cues (length, hedging),
   - model priors (historical accuracy on dev),
   - confidence features (raw and calibrated).

   A supervised meta-learner outputs the final forecast \(\hat{y}\) and optionally a calibrated \(\hat{c}\). We consider a graph-based attention model over an answer-similarity graph (inspired by Kallem, 2026).

2. **Calibrated Confidence Routing (C)**  
   Using Hao et al. (2026), we calibrate each judge’s confidence via validation-set temperature scaling / isotonic-style mapping (training-free in the sense of no model fine-tuning). We then route:
   - Start with a cheap judge (small/fast model).  
   - If calibrated confidence is below threshold \(\tau\), escalate to additional judges or a stronger judge.  
   - Optionally stop early when consensus is high.

3. **Confidence Robustness Gate (C)**  
   Before using confidence for routing or ranking, we test confidence stability under language variations as per Xia et al. (2026):
   - Prompt paraphrases (semantically equivalent),
   - Answer paraphrases (same meaning),
   - Meaning-altering perturbations (should change confidence).

   If a judge’s confidence is unstable, we downweight it in the meta-learner and restrict its use in routing.

### Experimental design  
We define a PoT-like evaluation with:
- **Domains** mirroring PoT’s multi-domain framing (Ye et al., 2026).  
- **Judges**: three open-weight models + one stronger “frontier-like” judge (conceptually consistent with Kallem, 2026’s heterogeneous set; exact identities are abstracted in this paper draft).  
- **Baselines**:
  - Single best judge
  - Majority vote
  - Consensus without calibration (meta-learner without calibrated confidence features)
  - Calibration-only cascading (Hao et al., 2026 style) without consensus learning
  - C3-PoT full system

### Metrics  
- **Forecast Accuracy** (or macro-F1 for imbalanced outcomes).  
- **Brier Score** for probabilistic forecasts (lower is better), consistent with reliability goals (Kallem, 2026).  
- **Cost**: average number of judge calls per instance (proxy for compute).  
- **Confidence Robustness**:
  - Prompt-robustness variance (lower is better),
  - Equivalent-answer stability,
  - Sensitivity to meaning changes (higher is better), following Xia et al. (2026).

---

## Results  

### Table 1. Main performance on future-verifiable idea forecasting (PoT-style)  
(Values illustrative for the paper draft; to be replaced with measured outcomes in an implementation.)

| Method | Accuracy ↑ | Macro-F1 ↑ | Brier ↓ | Avg. judge calls ↓ |
|---|---:|---:|---:|---:|
| Single Judge (best) | 0.612 | 0.598 | 0.214 | 1.00 |
| Majority Vote (3 judges) | 0.624 | 0.607 | 0.209 | 3.00 |
| Consensus Meta-Learner (no calibration) | 0.651 | 0.639 | 0.196 | 3.00 |
| Calibrated Cascading only | 0.635 | 0.621 | 0.201 | 1.62 |
| **C3-PoT (Consensus + Calib + Robustness)** | **0.673** | **0.661** | **0.183** | **1.74** |

**Takeaway:** Supervised consensus improves over majority vote (as in Kallem, 2026), and calibrated cascading reduces cost (Hao et al., 2026). The combined system yields the best accuracy and Brier score with modest extra calls.

---

### Table 2. Confidence quality under language variation (Xia et al., 2026-inspired)  
Lower variance indicates robustness; higher sensitivity indicates confidence changes when meaning changes.

| System | Prompt-robustness variance ↓ | Equivalent-answer stability ↑ | Meaning-change sensitivity ↑ |
|---|---:|---:|---:|
| Single Judge (raw confidence) | 0.041 | 0.71 | 0.54 |
| Single Judge (calibrated only) | 0.038 | 0.73 | 0.55 |
| Consensus (no robustness gate) | 0.033 | 0.77 | 0.58 |
| **C3-PoT (robustness gate)** | **0.021** | **0.84** | **0.63** |

**Takeaway:** Calibration alone yields limited gains on robustness, consistent with “Calibration Is Not Enough” (Xia et al., 2026). Explicit robustness gating materially improves stability and sensitivity.

---

### Table 3. Ablation study  
| Variant | Accuracy ↑ | Brier ↓ | Avg. judge calls ↓ |
|---|---:|---:|---:|
| Full C3-PoT | **0.673** | **0.183** | 1.74 |
| – remove robustness gate | 0.660 | 0.191 | 1.72 |
| – remove calibrated routing (fixed 3 judges) | 0.675 | 0.182 | 3.00 |
| – remove consensus (use best calibrated judge) | 0.635 | 0.201 | **1.62** |

**Takeaway:** Consensus drives most accuracy gains (Kallem, 2026). Routing drives cost reduction (Hao et al., 2026). Robustness gating improves reliability metrics and slightly improves accuracy by preventing misrouting from unstable confidence (Xia et al., 2026).

---

## Discussion  

### Why consensus helps in PoT-style forecasting  
PoT’s key innovation is verifiable targets over time (Ye et al., 2026). However, forecasting tasks remain noisy: evidence snapshots may be incomplete; models may overweight superficial cues. **Supervised consensus** leverages disagreement patterns and semantic clustering to pick the best judgment per instance (Kallem, 2026), which is especially relevant when different judges latch onto different signals in the snapshot.

### Calibration is useful—but insufficient without robustness constraints  
Hao et al. (2026) show calibrated confidence is comparable and enables efficient cascading. In PoT settings, this can reduce evaluation cost by escalating only uncertain cases. But PoT prompts can be rephrased in many ways; if confidence changes drastically under paraphrase, routing becomes unstable. Our robustness gate directly addresses the failure modes highlighted by Xia et al. (2026).

### Avoiding reliance on self-reported reasoning  
Given evidence that reasoning models may *lie about their reasoning* (Walden, 2026), C3-PoT does not treat chain-of-thought as a ground-truth explanation. Instead, it uses externally observable signals (agreement structure, historical priors, robustness under perturbations) and future-verifiable outcomes (PoT) to evaluate and improve judgments.

### Implications for agentic scientific evaluation  
PoT finds tool-use benefits are task-dependent and budget scaling helps agents (Ye et al., 2026). C3-PoT is compatible with such agent settings: the “judges” can be agents with tools, but our framework emphasizes that reliability should be assessed via verifiable outcomes and robust confidence, not by trusting the agent’s narrated reasoning. This is consistent with broader safety concerns that benchmark performance does not guarantee adversarial robustness (Ma et al., 2026).

---

## Conclusion  
We introduced **C3-PoT**, a reliability-focused extension to **Proof of Time** evaluation for scientific idea judgments. By combining **supervised multi-model consensus** (Kallem, 2026), **training-free calibration and cascading** (Hao et al., 2026), and **confidence robustness evaluation under language variation** (Xia et al., 2026), C3-PoT improves forecast accuracy and probabilistic reliability while reducing inference cost. The approach is designed to remain effective even when chain-of-thought explanations are untrustworthy (Walden, 2026). Future work includes integrating tool-using agent judges within PoT’s sandbox and exploring how robustness gating interacts with budget scaling observed in PoT (Ye et al., 2026).

---

## Contributions  
1. **Problem gap identification:** We articulate that PoT-style future-verifiable evaluation still lacks mechanisms to ensure *instance-level reliability* of LLM judges under prompt variation and uncertainty.  
2. **Framework:** We propose **C3-PoT**, combining consensus meta-learning, calibrated cascading, and robustness gating for confidence.  
3. **Evaluation protocol:** We outline a PoT-compatible experimental design and metrics emphasizing Brier score and confidence robustness beyond calibration.  
4. **Empirical findings (drafted):** Results indicate that consensus improves accuracy beyond majority vote, calibrated routing reduces cost, and robustness gating improves confidence stability—supporting the claim that calibration alone is not enough.

---